% Options for packages loaded elsewhere
\PassOptionsToPackage{unicode}{hyperref}
\PassOptionsToPackage{hyphens}{url}
\documentclass[
]{article}
\usepackage{xcolor}
\usepackage[margin=1in]{geometry}
\usepackage{amsmath,amssymb}
\setcounter{secnumdepth}{-\maxdimen} % remove section numbering
\usepackage{iftex}
\ifPDFTeX
  \usepackage[T1]{fontenc}
  \usepackage[utf8]{inputenc}
  \usepackage{textcomp} % provide euro and other symbols
\else % if luatex or xetex
  \usepackage{unicode-math} % this also loads fontspec
  \defaultfontfeatures{Scale=MatchLowercase}
  \defaultfontfeatures[\rmfamily]{Ligatures=TeX,Scale=1}
\fi
\usepackage{lmodern}
\ifPDFTeX\else
  % xetex/luatex font selection
\fi
% Use upquote if available, for straight quotes in verbatim environments
\IfFileExists{upquote.sty}{\usepackage{upquote}}{}
\IfFileExists{microtype.sty}{% use microtype if available
  \usepackage[]{microtype}
  \UseMicrotypeSet[protrusion]{basicmath} % disable protrusion for tt fonts
}{}
\makeatletter
\@ifundefined{KOMAClassName}{% if non-KOMA class
  \IfFileExists{parskip.sty}{%
    \usepackage{parskip}
  }{% else
    \setlength{\parindent}{0pt}
    \setlength{\parskip}{6pt plus 2pt minus 1pt}}
}{% if KOMA class
  \KOMAoptions{parskip=half}}
\makeatother
\usepackage{longtable,booktabs,array}
\usepackage{calc} % for calculating minipage widths
% Correct order of tables after \paragraph or \subparagraph
\usepackage{etoolbox}
\makeatletter
\patchcmd\longtable{\par}{\if@noskipsec\mbox{}\fi\par}{}{}
\makeatother
% Allow footnotes in longtable head/foot
\IfFileExists{footnotehyper.sty}{\usepackage{footnotehyper}}{\usepackage{footnote}}
\makesavenoteenv{longtable}
\usepackage{graphicx}
\makeatletter
\newsavebox\pandoc@box
\newcommand*\pandocbounded[1]{% scales image to fit in text height/width
  \sbox\pandoc@box{#1}%
  \Gscale@div\@tempa{\textheight}{\dimexpr\ht\pandoc@box+\dp\pandoc@box\relax}%
  \Gscale@div\@tempb{\linewidth}{\wd\pandoc@box}%
  \ifdim\@tempb\p@<\@tempa\p@\let\@tempa\@tempb\fi% select the smaller of both
  \ifdim\@tempa\p@<\p@\scalebox{\@tempa}{\usebox\pandoc@box}%
  \else\usebox{\pandoc@box}%
  \fi%
}
% Set default figure placement to htbp
\def\fps@figure{htbp}
\makeatother
\setlength{\emergencystretch}{3em} % prevent overfull lines
\providecommand{\tightlist}{%
  \setlength{\itemsep}{0pt}\setlength{\parskip}{0pt}}
\usepackage{bookmark}
\IfFileExists{xurl.sty}{\usepackage{xurl}}{} % add URL line breaks if available
\urlstyle{same}
\hypersetup{
  pdftitle={Анализ продаж магазина `Кофе и Книги'},
  pdfauthor={Товарный и временной анализ},
  hidelinks,
  pdfcreator={LaTeX via pandoc}}

\title{Анализ продаж магазина `Кофе и Книги'}
\author{Товарный и временной анализ}
\date{}

\begin{document}
\maketitle

{
\setcounter{tocdepth}{2}
\tableofcontents
}
\section{1. Общая
статистика}\label{ux43eux431ux449ux430ux44f-ux441ux442ux430ux442ux438ux441ux442ux438ux43aux430}

\begin{longtable}[]{@{}ll@{}}
\toprule\noalign{}
Показатель & Значение \\
\midrule\noalign{}
\endhead
\bottomrule\noalign{}
\endlastfoot
\textbf{Общая выручка} & 2,322,830 руб. \\
\textbf{Всего транзакций} & 2,000 \\
\textbf{Средний чек} & 1161.41 руб. \\
\textbf{Всего продано товаров} & 10,048 шт. \\
\end{longtable}

\begin{center}\rule{0.5\linewidth}{0.5pt}\end{center}

\section{2. Товарный
анализ}\label{ux442ux43eux432ux430ux440ux43dux44bux439-ux430ux43dux430ux43bux438ux437}

\subsection{2.1 Выручка по
категориям}\label{ux432ux44bux440ux443ux447ux43aux430-ux43fux43e-ux43aux430ux442ux435ux433ux43eux440ux438ux44fux43c}

\begin{longtable}[]{@{}lll@{}}
\caption{Выручка по категориям товаров}\tabularnewline
\toprule\noalign{}
Категория & Выручка (руб.) & Доля (\%) \\
\midrule\noalign{}
\endfirsthead
\toprule\noalign{}
Категория & Выручка (руб.) & Доля (\%) \\
\midrule\noalign{}
\endhead
\bottomrule\noalign{}
\endlastfoot
Книги & 950,450 & 40.9\% \\
Кофе & 688,810 & 29.7\% \\
Выпечка & 683,570 & 29.4\% \\
\end{longtable}

\pandocbounded{\includegraphics[keepaspectratio]{final_report_files/outputs/plots_A/revenue_by_category.png}}

\textbf{Вывод:} Основную выручку приносит категория \textbf{Книги} ---
40.9\% от общей выручки.

\subsection{2.2 Топ-5 товаров по
выручке}\label{ux442ux43eux43f-5-ux442ux43eux432ux430ux440ux43eux432-ux43fux43e-ux432ux44bux440ux443ux447ux43aux435}

\begin{longtable}[]{@{}lll@{}}
\caption{Топ-5 товаров по выручке}\tabularnewline
\toprule\noalign{}
Товар & Категория & Выручка (руб.) \\
\midrule\noalign{}
\endfirsthead
\toprule\noalign{}
Товар & Категория & Выручка (руб.) \\
\midrule\noalign{}
\endhead
\bottomrule\noalign{}
\endlastfoot
Мастер и Маргарита & Книги & 366,600 \\
Преступление и наказание & Книги & 299,000 \\
Война и мир & Книги & 284,850 \\
Сырник & Выпечка & 243,200 \\
Маффин & Выпечка & 188,850 \\
\end{longtable}

\pandocbounded{\includegraphics[keepaspectratio]{final_report_files/outputs/plots_A/top_products.png}}

\textbf{Вывод:} Топ-3 товара: Мастер и Маргарита, Преступление и
наказание, Война и мир.

\subsection{2.3 ABC-анализ}\label{abc-ux430ux43dux430ux43bux438ux437}

\begin{longtable}[]{@{}lrll@{}}
\caption{ABC-анализ товаров}\tabularnewline
\toprule\noalign{}
Класс ABC & Товаров & Выручка (руб.) & Доля (\%) \\
\midrule\noalign{}
\endfirsthead
\toprule\noalign{}
Класс ABC & Товаров & Выручка (руб.) & Доля (\%) \\
\midrule\noalign{}
\endhead
\bottomrule\noalign{}
\endlastfoot
A & 6 & 1,550,780 & 66.8\% \\
B & 3 & 463,170 & 19.9\% \\
C & 3 & 308,880 & 13.3\% \\
\end{longtable}

\pandocbounded{\includegraphics[keepaspectratio]{final_report_files/outputs/plots_A/abc_analysis.png}}

\textbf{Вывод:} Категория A приносит 70\% выручки, категория B --- 20\%,
категория C --- 10\%. \#\# 2.4 Цены на товары

\pandocbounded{\includegraphics[keepaspectratio]{final_report_files/outputs/plots_A/prices.png}}

\subsection{2.5 Выводы студента
А}\label{ux432ux44bux432ux43eux434ux44b-ux441ux442ux443ux434ux435ux43dux442ux430-ux430}

\begin{verbatim}

## 📊 Выводы по товарному анализу (Студент А)

### 1. Главный драйвер выручки
Категория **'Книги'** приносит 40.9% всей выручки.
- **Рекомендация:** расширить ассортимент в этой категории, выделить больше места на витрине.

### 2. Аутсайдеры (категория C)
3 товаров (категория C) приносят всего 10% выручки.
- **Рекомендация:** провести распродажу, сделать скидку 20-30% или убрать из ассортимента.

### 3. Ценовая политика
- Самый дорогой товар: **Мастер и Маргарита** (600 руб.)
- Самый дешёвый товар: **Печенье** (80 руб.)
- **Рекомендация:** добавить комбо-предложения, например 'Кофе + Выпечка = скидка 15%'.

### 4. Топ-5 товаров
Мастер и Маргарита, Преступление и наказание, Война и мир, Сырник, Маффин
- **Рекомендация:** всегда держать эти товары в наличии, сделать отдельную витрину.

### 5. ABC-анализ
- Товары категории A: 6 штук (70% выручки)
- Товары категории B: 3 штук (20% выручки)
- Товары категории C: 3 штук (10% выручки)
\end{verbatim}

\begin{center}\rule{0.5\linewidth}{0.5pt}\end{center}

\section{3. Временной
анализ}\label{ux432ux440ux435ux43cux435ux43dux43dux43eux439-ux430ux43dux430ux43bux438ux437}

\subsection{3.1 Выручка по дням
недели}\label{ux432ux44bux440ux443ux447ux43aux430-ux43fux43e-ux434ux43dux44fux43c-ux43dux435ux434ux435ux43bux438}

\begin{longtable}[]{@{}llrl@{}}
\caption{Выручка по дням недели}\tabularnewline
\toprule\noalign{}
День & Выручка (руб.) & Транзакций & Средний чек (руб.) \\
\midrule\noalign{}
\endfirsthead
\toprule\noalign{}
День & Выручка (руб.) & Транзакций & Средний чек (руб.) \\
\midrule\noalign{}
\endhead
\bottomrule\noalign{}
\endlastfoot
вс & 304,660 & 256 & 1,190.1 \\
пн & 298,520 & 254 & 1,175.3 \\
вт & 346,850 & 303 & 1,144.7 \\
ср & 350,370 & 307 & 1,141.3 \\
чт & 334,810 & 303 & 1,105.0 \\
пт & 337,880 & 280 & 1,206.7 \\
сб & 349,740 & 297 & 1,177.6 \\
\end{longtable}

{}

\textbf{Вывод:} Пик продаж --- \textbf{ср}.

\subsection{3.2 Средний чек по дням
недели}\label{ux441ux440ux435ux434ux43dux438ux439-ux447ux435ux43a-ux43fux43e-ux434ux43dux44fux43c-ux43dux435ux434ux435ux43bux438}

\pandocbounded{\includegraphics[keepaspectratio]{final_report_files/outputs/plots_B/avg_check_by_weekday.png}}

\subsection{3.3 Будни vs
выходные}\label{ux431ux443ux434ux43dux438-vs-ux432ux44bux445ux43eux434ux43dux44bux435}

\begin{longtable}[]{@{}llrl@{}}
\caption{Сравнение будней и выходных}\tabularnewline
\toprule\noalign{}
Тип & Выручка (руб.) & Транзакций & Средний чек (руб.) \\
\midrule\noalign{}
\endfirsthead
\toprule\noalign{}
Тип & Выручка (руб.) & Транзакций & Средний чек (руб.) \\
\midrule\noalign{}
\endhead
\bottomrule\noalign{}
\endlastfoot
Будни & 1,668,430 & 1447 & 1,153 \\
Выходные & 654,400 & 553 & 1,183 \\
\end{longtable}

\pandocbounded{\includegraphics[keepaspectratio]{final_report_files/outputs/plots_B/weekend_vs_weekday.png}}

\subsection{3.4 Динамика
продаж}\label{ux434ux438ux43dux430ux43cux438ux43aux430-ux43fux440ux43eux434ux430ux436}

\pandocbounded{\includegraphics[keepaspectratio]{final_report_files/outputs/plots_B/daily_revenue_trend.png}}

\subsection{3.5 Выводы студента
Б}\label{ux432ux44bux432ux43eux434ux44b-ux441ux442ux443ux434ux435ux43dux442ux430-ux431}

\begin{verbatim}
Выводы по временному анализу (Студент Б)

1. Самый прибыльный день
Среда показала максимальную выручку 350 370 рублей.
Понедельник - самый слабый день с выручкой 298 520 рублей.
Разница составляет 51 850 рублей (17.4%).
Рекомендация: проводить акции в понедельник для выравнивания продаж,
обеспечивать максимальный запас товаров (как книг, так и кофе) к среде.

2. Средний чек по дням недели
Самый высокий средний чек в пятницу - 1207 рублей.
Самый низкий в четверг - 1105 рублей.
Разница составляет 102 рублей (9.2%).
Рекомендация: в четверг запустить акцию "Книга + кофе = скидка"
или предлагать десерт/бестселлер к основному заказу для повышения чека.
Пятница показывает отличные результаты, достаточно поддерживать текущий уровень.

3. Сравнение будней и выходных
Средний чек в будни - 1153 рублей, в выходные - 1183 рублей.
Разница составляет всего 30 рублей (2.6%), что статистически незначительно.
Рекомендация: не требуется отдельной стратегии для выходных.
В будни можно тестировать акции на книги (когда посетителей меньше),
в выходные сохранять стандартное меню и ассортимент.

4. Итоговые рекомендации
Усилить понедельник акциями (самый слабый день по выручке).
Повышать средний чек в четверг через комбо-предложения "Книга + кофе".
Сохранить базовую стратегию для выходных - отличие от будней минимально.
Обеспечить достаточный запас товаров к среде - дню максимальной выручки.
В пятницу (максимальный чек) можно предлагать новинки книг и сезонные напитки.
\end{verbatim}

\begin{center}\rule{0.5\linewidth}{0.5pt}\end{center}

\section{4. Общие
рекомендации}\label{ux43eux431ux449ux438ux435-ux440ux435ux43aux43eux43cux435ux43dux434ux430ux446ux438ux438}

\subsubsection{Ассортиментная
стратегия}\label{ux430ux441ux441ux43eux440ux442ux438ux43cux435ux43dux442ux43dux430ux44f-ux441ux442ux440ux430ux442ux435ux433ux438ux44f}

\begin{itemize}
\tightlist
\item
  Укрепить категорию \textbf{Книги}
\item
  Провести распродажу товаров категории C
\end{itemize}

\subsubsection{Временная
стратегия}\label{ux432ux440ux435ux43cux435ux43dux43dux430ux44f-ux441ux442ux440ux430ux442ux435ux433ux438ux44f}

\begin{itemize}
\tightlist
\item
  Увеличить запасы к пятнице и субботе
\item
  Запустить акцию ``Счастливый вторник''
\end{itemize}

\subsubsection{Маркетинговые
активности}\label{ux43cux430ux440ux43aux435ux442ux438ux43dux433ux43eux432ux44bux435-ux430ux43aux442ux438ux432ux43dux43eux441ux442ux438}

\begin{enumerate}
\def\labelenumi{\arabic{enumi}.}
\tightlist
\item
  По вторникам --- скидки
\item
  В выходные --- премиум-меню
\item
  Для топ-товаров --- отдельная витрина
\end{enumerate}

\begin{center}\rule{0.5\linewidth}{0.5pt}\end{center}

\section{5.
Заключение}\label{ux437ux430ux43aux43bux44eux447ux435ux43dux438ux435}

\begin{longtable}[]{@{}llll@{}}
\toprule\noalign{}
Студент & Задачи & Графиков & Выводов \\
\midrule\noalign{}
\endhead
\bottomrule\noalign{}
\endlastfoot
\textbf{А} & Товарный анализ & 4 & 4 \\
\textbf{Б} & Временной анализ & 4 & 4 \\
\end{longtable}

\textbf{Дата защиты:} 2026-06-09

\begin{center}\rule{0.5\linewidth}{0.5pt}\end{center}

\end{document}
